\documentclass[12pt,a4paper]{report}

\usepackage[T1]{fontenc}
\usepackage[utf8]{inputenc}
\usepackage[czech,english]{babel}
\usepackage{lmodern}
\usepackage{graphicx}
\usepackage{hyperref}
\usepackage{natbib}
\usepackage[final]{microtype}
\usepackage{xspace}
\usepackage[noabbrev,capitalize]{cleveref}

\setlength\textwidth{145mm}
\setlength\textheight{247mm}
\setlength\oddsidemargin{15mm}
\setlength\evensidemargin{15mm}
\setlength\topmargin{0mm}
\setlength\headsep{0mm}
\setlength\headheight{0mm}
\setlength{\parskip}{0.35em}
\setlength{\parindent}{0pt}
\setcounter{tocdepth}{4}
\raggedbottom

\hypersetup{
  unicode=true,
  breaklinks=true,
  bookmarksdepth=0,
  pdftitle={Abstract of Doctoral Thesis},
  pdfauthor={Petra Barančíková}
}

\newcommand{\ThesisTitle}{Evaluating Metrics for Machine Translation Evaluation: From Paraphrases to Semantic Embeddings}
\newcommand{\ThesisTitleCzech}{Hodnocení metrik pro hodnocení strojového překladu: od parafrází k sémantickým embeddingům}
\newcommand{\ThesisAuthor}{Mgr. Petra Barančíková}
\newcommand{\Department}{Institute of Formal and Applied Linguistics}
\newcommand{\DepartmentCzech}{Ústav formální a aplikované lingvistiky}
\newcommand{\Supervisor}{doc. RNDr. Markéta Lopatková, Ph.D.}
\newcommand{\Consultant}{prof. RNDr. Ondřej Bojar, Ph.D.}
\newcommand{\StudyProgramme}{Computational Linguistics}
\newcommand{\StudyProgrammeCzech}{Počítačová lingvistika}
\newcommand{\StudyBranch}{Mathematical Linguistics}
\newcommand{\StudyBranchCzech}{Matematická lingvistika}
\newcommand{\YearSubmitted}{2026}
\newcommand{\FacultyEN}{Faculty of Mathematics and Physics}
\newcommand{\FacultyCZ}{Matematicko-fyzikální fakulta}
\newcommand{\UniversityEN}{Charles University}
\newcommand{\UniversityCZ}{Univerzita Karlova}
\newcommand{\FacultyAddress}{Malostranské nám. 25, 118 00 Praha 1}
\newcommand{\StudyDepartmentAddress}{Ke Karlovu 3, 121 16 Praha 2}
\newcommand{\DefenceDateEN}{June 24, 2026}
\newcommand{\DefenceDateCZ}{22. září 2026}
\newcommand{\DefenceRoom}{Malá aula}
\newcommand{\OpponentsEN}{doc. Mgr. Pavel Rychlý, Ph.D.\\ doc. Jernej Vičič, Ph.D.}
\newcommand{\OpponentsCZ}{doc. Mgr. Pavel Rychlý, Ph.D.\\ doc. Jernej Vičič, Ph.D.}

\newcommand{\bleu}{\textsc{BLEU}\xspace}
\newcommand{\ter}{\textsc{TER}\xspace}
\newcommand{\meteor}{\textsc{Meteor}\xspace}
\newcommand{\mn}{\textsc{Meteor-NEXT}\xspace}
\newcommand{\comet}{\textsc{COMET}\xspace}
\newcommand{\bleurt}{\textsc{BLEURT}\xspace}
\newcommand{\parmesan}{\textsc{Parmesan}\xspace}
\newcommand{\costra}{\textsc{Costra}\xspace}
\newcommand{\costraminus}{\textsc{Costra--}\xspace}
\newcommand{\mpt}{the \textsc{Meteor} paraphrase tables\xspace}
  
\newcommand{\blankpage}{\newpage\thispagestyle{empty}\null\newpage}
\newcommand{\chapwithtoc}[1]{\chapter*{#1}\addcontentsline{toc}{chapter}{#1}}
\newcommand{\frontmattertitle}[1]{{\fontsize{16}{18}\selectfont\bfseries #1}}
\newcommand{\frontmatterline}[1]{#1\par}
\newcommand{\coverauthor}[1]{{\fontsize{16}{18}\selectfont #1}}
\newcommand{\covertitle}[1]{{\fontsize{18}{22}\selectfont\bfseries\parbox{0.74\textwidth}{\centering #1}}}
\newenvironment{detailpage}{%
  \thispagestyle{empty}%
  \vspace*{12mm}%
  \small
  \setlength{\parskip}{0.3em}%
  \setlength{\parindent}{0pt}%
}{%
}
\newcommand{\detailrow}[2]{%
  \noindent\begin{tabular}{@{}p{35mm}@{\hspace{3mm}}p{102mm}@{}}%
    #1 & \begin{minipage}[t]{102mm}#2\end{minipage}
  \end{tabular}\par%
}

\begin{document}
\selectlanguage{english}
\pagestyle{empty}
\hypersetup{pageanchor=false}

\begin{center}
\vspace*{10mm}
\includegraphics[width=120mm]{../img/logo-en.pdf}

\vspace{38mm}
\frontmattertitle{ABSTRACT OF DOCTORAL THESIS}

\vspace{14mm}
\coverauthor{\ThesisAuthor}

\vspace{18mm}
\covertitle{\ThesisTitle}

\vspace{20mm}
{\Department}

\vspace{18mm}
{\renewcommand{\arraystretch}{1.15}
\begin{tabular}{rl}
Supervisor: & \Supervisor \\
Consultant: & \Consultant \\
Study programme: & \StudyProgramme \\
Study branch: & \StudyBranch
\end{tabular}
}

\vfill
Prague \YearSubmitted
\end{center}

\blankpage

\selectlanguage{english}
\begin{detailpage}
The results of this thesis were achieved in the period of doctoral study at the \FacultyEN, \UniversityEN.

\vspace{1.2em}
\detailrow{Student:}{\ThesisAuthor}
\detailrow{Supervisor:}{\Supervisor\\\Department\\\FacultyEN\\\UniversityEN\\\FacultyAddress}
\detailrow{Consultant:}{\Consultant\\\Department\\\FacultyEN\\\UniversityEN\\\FacultyAddress}

\vspace{0.7em}
\detailrow{Department:}{\Department\\\FacultyEN\\\UniversityEN\\\FacultyAddress}

\vspace{0.7em}
\detailrow{Opponents:}{doc. Mgr. Pavel Rychlý, Ph.D.\\Masarykova univerzita\\[0.4em]doc. Jernej Vičič, Ph.D.\\University of Primorska}

\vspace{1.2em}
The thesis defence will take place on \DefenceDateEN\ before the committee for thesis defences in Mathematical Linguistics at the \FacultyEN, \UniversityEN, \FacultyAddress, room \DefenceRoom.

\vspace{1.2em}
The thesis can be viewed at the Study Department of Doctoral Studies of the \FacultyEN, \UniversityEN, \StudyDepartmentAddress.
\end{detailpage}

\newpage

\selectlanguage{czech}
\begin{center}
\vspace*{10mm}
\includegraphics[width=120mm]{../img/logo-cs.pdf}

\vspace{38mm}
\frontmattertitle{AUTOREFERÁT DISERTAČNÍ PRÁCE}

\vspace{14mm}
\coverauthor{\ThesisAuthor}

\vspace{18mm}
\covertitle{\ThesisTitleCzech}

\vspace{20mm}
{\DepartmentCzech}

\vspace{18mm}
{\renewcommand{\arraystretch}{1.15}
\begin{tabular}{rl}
Školitelka: & \Supervisor \\
Konzultant: & \Consultant \\
Studijní program: & \StudyProgrammeCzech \\
Studijní obor: & \StudyBranchCzech
\end{tabular}
}

\vfill
Praha \YearSubmitted
\end{center}

\blankpage

\selectlanguage{czech}
\begin{detailpage}
Výsledky uvedené v této disertační práci byly získány během doktorského studia na \FacultyCZ, \UniversityCZ.

\vspace{1.2em}
\detailrow{Doktorandka:}{\ThesisAuthor}
\detailrow{Školitelka:}{\Supervisor\\\DepartmentCzech\\\FacultyCZ\\\UniversityCZ\\\FacultyAddress}
\detailrow{Konzultant:}{\Consultant\\\DepartmentCzech\\\FacultyCZ\\\UniversityCZ\\\FacultyAddress}

\vspace{0.7em}
\detailrow{Pracoviště:}{\DepartmentCzech\\\FacultyCZ\\\UniversityCZ\\\FacultyAddress}

\vspace{0.7em}
\detailrow{Oponenti:}{doc. Mgr. Pavel Rychlý, Ph.D.\\Masarykova univerzita\\[0.4em]doc. Jernej Vičič, Ph.D.\\Univerza na Primorskem}

\vspace{1.2em}
Obhajoba disertační práce se koná dne \DefenceDateCZ\ před komisí pro obhajoby disertačních prací v oboru \StudyBranchCzech\ na \FacultyCZ, \UniversityCZ, \FacultyAddress, místnost \DefenceRoom.

\vspace{1.2em}
S disertační prací je možné se seznámit na oddělení doktorského studia \FacultyCZ, \UniversityCZ, \StudyDepartmentAddress.
\end{detailpage}



\newpage
\selectlanguage{english}
\tableofcontents
\clearpage
\hypersetup{pageanchor=true}
\pagestyle{plain}
\pagenumbering{arabic}

\chapwithtoc{Introduction}

During the last decade, machine translation (MT) systems have changed significantly
progressing from phrase-based statistical systems to neural models and, more recently, 
to large language models. 
Although translation quality has improved dramatically, reliable evaluation of translation 
outputs remains challenging. Reliable evaluation is still very important: 
it is needed to compare systems, guide model development, interpret shared-task results, 
and make research claims comparable across time.

Human evaluation remains the gold standard for assessing translation quality. However,
it is expensive, time-consuming, subjective, and difficult to reproduce at scale.
Consequently, automatic metrics are important component of MT research. 

The main challenge for machine translation is that a single source sentence can be 
translated correctly in many different ways. A useful metric must decide whether a 
hypothesis and a reference translation, i.e. sentence translated by professional human
translator, express the same meaning even when they differ lexically, morphologically, 
or syntactically.

Rich morphology and relatively free word order of the Czech language allow many surface 
variants that preserve the same meaning. Traditional string-based metrics such as \bleu 
can penalize correct translations simply because they differ from the reference wording 
(see \Cref{example_of_BLEU_malfunction}). Modern neural metrics reduce this limitation 
by relying on learned semantic representations, but they introduce new questions: what 
do sentence embeddings actually encode, can they reliably encode the meaning of sentence?

 \begin{figure}[tb]
\centering
\small
\resizebox{\linewidth}{!}{%
\begin{tabular}{ll}
\toprule
\textbf{Source} &
\textit{Banks are testing payment by mobile telephone} \\
\midrule
\textbf{Hyp.} &
\begin{tabular}[t]{@{}llllll@{}}
\foreign{Banky} & \foreign{zkoušejí} & \foreign{platbu} & \foreign{pomocí} & \foreign{mobilního} & \foreign{telefonu} \\
\trans{Banks}   & \trans{are testing}& \trans{payment}  & \trans{by}       & \trans{mobile} & \trans{phone}
\end{tabular} \\
&
\trans{Banks are testing payment by mobile phone} \\
\midrule
\textbf{Ref.} &
\begin{tabular}[t]{@{}llll@{}}
\foreign{Banky} & \foreign{testují} & \foreign{placení} & \foreign{mobilem} \\
\trans{Banks}   & \trans{are testing} & \trans{paying}  & \trans{by mobile phone} \\
\end{tabular} \\
&
\trans{Banks are testing paying by mobile phone} \\
\bottomrule
\end{tabular}
}
\caption{In this example from  Workshop on Machine Translation (WMT) 2012 \citep{wmt12}, despite the hypothesis being grammatically correct and having the same meaning as the source sentence, there is only a single overlapping unigram between the hypothesis and the reference translation. As a result, the hypothesis would receive a very low score when evaluated using n-gram overlap-based automatic evaluation metrics.}
\label{example_of_BLEU_malfunction}
\end{figure}

This dissertation follows both sides of this development. The first part revisits the 
traditional setting and asks whether string-based evaluation can be improved by 
automatically paraphrasing reference translations. The second part studies sentence 
embeddings as a modern semantic alternative, examining their geometric structure, their 
behavior on Czech semantic transformations, and their performance inside neural MT 
evaluation metrics.

Although the underlying technologies differ substantially, both parts address the same 
research question:  how automatic evaluation can better capture semantic equivalence 
between translations. More specifically, the thesis addresses four research questions: 

\begin{itemize}
\item Can the reliability of traditional string based methods be improved through 
      targeted paraphrasing of reference sentences?
\item Do sentence embeddings exhibit systematic geometric structure corresponding to semantic transformations?
\item How well do modern sentence embedding models capture complex semantic transformations in Czech?
\item Does strong perfomance on intrinsic semantic benchmarks imply strong performance as
       a machine translation evaluation metric?
\end{itemize}

\chapter{Machine Translation Evaluation}

The thesis first establishes the evaluation setting in which the experiments are interpreted. 
We present a detailed survey of machine translation evaluation methods from the last twenty 
years at the Workshop on Machine Translation (WMT) shared tasks from adequacy and fluency 
judgments through relative ranking to Direct Assessment, MQM, and Error Span Annotation.
These protocols differ in cost, granularity, and reproducibility, but they share the same 
goal: to estimate how much a translation preserves the meaning of the source. They are also 
important because their scores serve as training or validation targets for many automatic metrics. 

Automatic evaluation metrics attempt to approximate human judgments efficiently and reproducibly. 
Traditional metrics such as \bleu \citep{papineni-etal-2002-bleu}, \ter, and character-level 
metrics measure surface overlap between a system hypothesis and one or more human references. 
Their advantages are practical: they are fast, deterministic, transparent, and comparable with 
a large body of previous work. Their weakness is conceptual: lexical overlap is only an imperfect 
proxy for translation quality. Correct translations may use different words, inflectional forms, 
or word order than the reference.

This limitation motivated metrics that use linguistic resources. Metrics such as \meteor 
and \mn add stemming, synonym matching, and paraphrase tables, making the comparison more 
tolerant of legitimate variation. However, the added resources also introduce noise and 
language dependence. In the Czech setting, automatically created paraphrase tables contain 
many pairs that are not valid paraphrases in context. When these pairs are used directly 
inside a metric, the metric can reward grammatical errors, inappropriate substitutions, or 
partially untranslated output.

Recent neural metrics, including \bleurt \citep{sellam2020bleurt} and \comet 
\citep{rei-etal-2020-comet}, shift the problem from explicit string matching to learned 
semantic representations. They often correlate better with human judgments, but they are 
less transparent and depend on training data, model architecture, and evaluation protocol. 
At the same time, traditional metrics remain widely reported in MT research 
\citep{marie-etal-2021-scientific}. For this reason, the thesis studies both paradigms: 
explicit paraphrasing for string-based metrics and sentence embeddings for neural semantic 
evaluation.

\chapter{Targeted Sentence Paraphrasing for MT Evaluation}

The first experimental part investigates targeted paraphrasing of reference translations. 
The motivation is simple: if a metric penalizes a valid hypothesis because it is lexically 
different from the reference, the evaluation can be improved by generating additional 
reference variants that are closer to the hypothesis while preserving the original meaning. 
This approach keeps the metric itself mostly unchanged and instead improves the material 
against which the hypothesis is compared.

\paragraph{Simple substitution.}
The first method is deliberately conservative. It identifies candidate paraphrase pairs 
between a hypothesis and its reference, validates them against lexical and phrasal paraphrase 
resources, substitutes selected expressions into the reference, and optionally applies Depfix 
to repair grammatical errors introduced by the substitution. The method is inspired by earlier 
work on paraphrasing for automatic evaluation \citep{kauchak}, but it is adapted to Czech
morphology and to the limited reliability of Czech paraphrase resources.

This work demonstrated that evaluation quality can be improved without modifying \bleu 
itself, simply by enriching the reference space. The strongest improvements come from 
one-word substitutions selected from higher-quality resources and followed by grammatical 
correction. Multi-word paraphrases increase coverage, but they are more vulnerable to 
noise. The main lesson is that paraphrase precision and grammatical well-formedness 
matter more than the complexity of the generation method.

\paragraph{\parmesan.}
The next experiment asks whether paraphrased references are still useful for metrics 
that already include paraphrase matching internally, especially \mn. Several variants 
of \mn are compared with progressively filtered paraphrase tables and with three reference 
types: original references, paraphrased references, and Depfix-corrected paraphrased 
references. The best configuration removes internal paraphrase matching from \mn and 
evaluates Depfix-corrected paraphrased references using exact matching.

This leads to the \parmesan metric, which shows that paraphrases are more effective 
as alternative references than as internal matching rules. The result is conceptually 
important: paraphrases help when they become grammatical alternative references, but
noisy paraphrase tables can harm evaluation when they are used directly as metric-internal 
matching rules. The method therefore separates two tasks that are often conflated: 
generating plausible reference variants and scoring the final hypothesis-reference match.


\paragraph{Paraphrasing via machine translation.}
The thesis also explores whether more advanced methods of paraphrasing can improve on 
the simple substitution method. Paraphrasing and translation are closely related tasks, 
they both aim to preserve meaning of a sentence while changing its surface form, however
there are much more resources and tools available for translation than for paraphrasing.
Therefore, we adapted existing machine translation systems for monolingual paraphrasing. 

First, we adapt Moses, a phrase-based statistical MT system, for monolingual Czech-to-Czech 
paraphrasing. Czech paraphrase tables are used as translation tables, and a new decoder 
feature guides generation toward paraphrases that overlap with the evaluated hypothesis. 
Further, we experiment with adapting Treex, a deep-syntactic transfer-based MT system, 
for monolingual targeted paraphrasing.

These systems are more expressive, but they do not consistently outperform the simpler 
method. Their outputs can introduce broader structural changes and additional errors, 
making the resulting references less controlled and sometimes less useful for evaluation. 
This finding is one of the main empirical results of the first part: for MT evaluation, 
a simple, precise, grammatically corrected paraphrase can be more valuable than a 
sophisticated generator.

\paragraph{Findings.}
Overall, the paraphrasing experiments show that reference expansion can improve 
traditional metrics for Czech MT evaluation, but only when paraphrases are targeted, 
precise, and grammatically sound. Larger paraphrase resources are not automatically
better because noise can outweigh coverage. The results remain relevant even in the 
era of neural metrics: string-based metrics are still widely reported, and the 
paraphrasing experiments make explicit where surface matching fails.

\chapter{Sentence Embedding Evaluation}

In the second part, we focus on sentence embeddings evaluation. Embedding-based metrics 
represent sentences as points in a continuous vector space, where geometric similarity 
is expected to correspond to semantic similarity. This decreases the need for explicit
paraphrasing, but it introduces a new question: what semantic information is actually 
encoded in sentence embeddings, and is it encoded in a way that is useful for MT evaluation?

\paragraph{Linear relations in embedding spaces.}
In the first section of the second part, we investigate whether sentence embeddings 
exhibit consistent linear structures analogous to word embeddings. As \cite{word2vec} 
demonstrated, word embeddings capture semantic and syntactic relationships geometrically, 
which they evaluated through tests like adjective declension (easy'' $\rightarrow$ easier'' 
→ easiest''), verb tense changes (walking'' → walk''), and capital-country matching (Athens'' 
→ ``Greece'').

\begin{figure}[t!]
  \centering
  \includegraphics{spaces.png}
\caption{An illustration reprinted from \cite{barancikova-bojar-linear}. The upper plot 
represents a continuous multi-dimensional vector space where individual sentences are 
represented as points (the ``space of sentences''). Pairs of related sentences are 
connected by arrows, with different dashing styles indicating various relationship types.\\
The lower plot illustrates the corresponding ``space of operations.'' Here, a vector 
difference is shown, with all arrows moved to start at a common origin. \\
The underlying hypothesis is that consistent semantic changes (e.g., all vector 
transformations extracted from sentence pairs differing in the speed of travel ``running 
instead of walking'') should be represented close to each other in the space of operations, 
i.e., form a more or less compact \emph{cluster}.}
\label{space-of-sentences}
\end{figure}

Our objective is to determine whether \emph{whole sentence} embeddings exhibit the same linear 
behavior, looking for an analogous behavior to that observed with word embeddings. This idea is
illustrated in \cref{space-of-sentences}. Specifically, we examine whether semantic shifts 
between sentences can be modeled as transformation vectors. Understanding whether sentence
embeddings exhibit consistent linear structure is particularly important for embedding-based 
evaluation metrics, which implicitly rely on geometric properties such as vector similarity 
and compositionality.

The results show that some sentence embedding models do contain coherent geometric structure. 
Certain lexical and semantic changes form recognizable clusters, indicating that the models 
encode more than arbitrary surface similarity. At the same time, this structure depends 
strongly on the training objective. Architecture alone is not sufficient: untrained or 
poorly suited models do not show the same behavior. Negation and more complex logical changes 
remain difficult across models.

\paragraph{The \costra dataset.}
The major contribution of this thesis is \costra, a Czech benchmark for evaluating 
sentence embeddings \citep{costra10,costra11}. It consists of 7 thousand 
manually created sentences and 16 thousand pairwise human similarity judgments 
for evaluation. The dataset was designed to test whether sentence embeddings 
capture complex semantic relations beyond surface similarity. The sentences
are selected to create a real challenge for the embeddings, e.g., what is closer to original
sentence - a distinct paraphrase or sentence that is similar on the surface but different in 
meaning?  Other categories include temporal changes, modality, opposition, generalization or
style changes.

Experiments on \costra show that current sentence encoders still rely heavily on surface 
similarity. Even strong models struggle when lexical overlap is high but meaning changes 
substantially, and they often fail to capture ordinal semantic progressions. Czech-specific 
and multilingual models each have strengths, but no tested model fully solves the diagnostic 
tasks. The \costra results therefore reveal limitations that are often hidden by standard 
semantic textual similarity benchmarks.

\paragraph{Intrinsic vs. extrinsic evaluation.}

In the previous section, we demonstrated through intrinsic evaluation that sentence 
embeddings often rely on surface-level similarities rather than capturing complex 
linguistic phenomena. However, the primary focus of this thesis is extrinsic 
evaluation --- machine translation evaluation (MTE). 

In contrast to the simple mathematical computations used in intrinsic evaluation, 
extrinsic evaluation for MTE requires costly model training. This leads to a natural 
question: can intrinsic evaluation serve as a predictor for downstream performance? 
Semantic Textual Similarity (STS) is closely related to MTE, as both tasks assess 
whether two sentences convey the same meaning despite surface differences. If intrinsic 
benchmarks correlated well with extrinsic performance, they could be used as a quick 
proxy for selecting promising embedding models. 

This question is particularly relevant given the growing number of available embedding 
models. Multilingual encoders such as \labse and \xlmr provide strong cross-lingual 
representations, while language-specific models aim to better capture properties of 
individual languages. For morphologically rich languages such as Czech, it is not 
obvious which approach is more suitable for MT evaluation. This raises another question: 
do Czech-specific encoders outperform multilingual models in evaluation of Czech translations?

In the final experiments, we compare intrinsic evaluation of Czech sentence embeddings 
with their downstream usefulness in MT evaluation \citep{barancikova-bojar-2025-intrinsic}. 
Several Czech and multilingual encoders are evaluated on \costra and Czech STS datasets. 
The same encoders are then used inside \comet-based metrics trained for machine translation 
evaluation and quality estimation on WMT English-to-Czech data.

The results reveal a striking mismatch. Models that perform best on intrinsic semantic 
benchmarks do not necessarily become the best MT evaluation metrics after fine-tuning. 
Conversely, some models with weak intrinsic performance become highly competitive 
inside \comet-based metrics. Task-specific fine-tuning changes the model ranking 
substantially, which challenges the assumption that intrinsic semantic quality is a 
reliable shortcut for choosing encoders for MT evaluation.

The comparison also shows that reference-free quality estimation can become competitive 
with reference-based evaluation in this setup. This agrees with recent developments in
MT evaluation, where source-aware and reference-free neural metrics often perform on par 
with reference-based metrics. The thesis therefore concludes that embeddings must be selected 
and validated with respect to the target task, not only with respect to standalone semantic 
benchmarks.

\chapwithtoc{Conclusion}

This thesis examines two ways of addressing the same fundamental problem: automatic 
metrics must evaluate meaning, but meaning is not tied to a single surface form. The 
first part uses explicit paraphrases to make traditional metrics more tolerant of valid 
translation variants. The second part analyzes sentence embeddings as implicit semantic 
representations used by modern neural metrics.

Across both parts, the same conclusion emerges: the most complex method is not necessarily 
the most effective one. Simple lexical paraphrases, when selected carefully and corrected 
grammatically, improve evaluation more reliably than more ambitious paraphrase generators. 
Similarly, larger or intrinsically stronger embedding models do not automatically produce 
better MT evaluation metrics after fine-tuning. What matters is alignment between the method, 
the data, and the specific notion of semantic equivalence required by the evaluation task.

The dissertation contributes practical resources and methodological insights: targeted 
paraphrasing methods for Czech MT evaluation, the \parmesan metric, an analysis of linear 
structure in sentence embeddings, the \costra benchmark for Czech sentence embeddings, 
and an empirical comparison of intrinsic and extrinsic embedding evaluation. Together, 
these contributions document the transition from surface-based evaluation toward semantic 
evaluation while showing that semantic evaluation remains difficult, language-sensitive, 
and task-dependent.

This thesis explores two complementary approaches to automatic machine translation 
evaluation to address the challenges of measuring semantic equivalence: targeted 
paraphrasing for traditional string-based metrics and the analysis of modern sentence 
embeddings. 

----
This thesis examines two complementary approaches to automatic machine translation evaluation. The first part focuses on targeted paraphrasing of reference translations as a method to improve the reliability of traditional string-based evaluation metrics. The second part examines the behavior and limitations of modern sentence embedding models with respect to machine translation evaluation. Together, these two approaches highlight the difficulty of measuring semantic equivalence in machine translation.

In addition to the experimental work, the thesis provides a structured overview of machine translation evaluation methods, with a particular focus on human evaluation in the WMT shared tasks. We trace the development from early adequacy and fluency judgments through relative ranking to Direct Assessment, MQM, and more recent approaches such as ESA. This overview clarifies what current evaluation protocols actually measure, as well as their practical limitations, including subjectivity, cost, and limited reproducibility. It also provides an essential context for interpreting automatic metrics, which are typically trained to approximate these human judgments.

Building upon this theoretical foundation, the main contributions of this thesis can be summarized as follows. First, we demonstrate that the reliability of string-based metrics such as \bleu can be systematically improved --- specifically in the context of a morphologically rich language like Czech --- by creating artificial reference sentences closer to machine translation outputs. We show that for Czech, simple substitution-based paraphrasing methods consistently outperform more complex approaches based on statistical or rule-based machine translation systems. This finding suggests that controlled lexical variation is more effective than generating complete paraphrases that differ more significantly in structure. At the same time, we show that large, automatically created paraphrase resources can reduce evaluation quality by introducing noise. Filtering and moving paraphrase processing from evaluation to preprocessing improves results. Finally, we establish that grammatical correctness plays a crucial role: linguistically well-formed paraphrases consistently improve evaluation reliability.

Second, we provide a detailed analysis of sentence embeddings as a semantic evaluation tool. We analyze the geometric properties of both English and multilingual sentence embedding spaces and show that certain models exhibit a clear and interpretable geometric structure, in which simple linguistic transformations correspond to consistent directional shifts. This structure is not a trivial byproduct of architecture, but emerges from the training objective, particularly in models trained with semantic or contrastive signals.

To support further examination of embedding properties, we introduce the \costra dataset of complex Czech sentence transformations, addressing a gap in existing evaluation resources. Our experiments show that even strong, state-of-the-art embedding models continue to rely heavily on surface-level similarity and often fail to capture more complex semantic relations. We also find that strong performance on intrinsic evaluation tasks does not reliably predict performance in machine translation evaluation, despite both being designed to assess semantic similarity.

Taken together, the results lead to a consistent conclusion: a more complex model does not automatically result in better evaluation performance. Simple string-based metrics, when supported by carefully constructed and grammatically sound paraphrases, can perform as well as or better than more complex approaches. Similarly, larger and more general sentence embedding models do not necessarily outperform smaller or more specialized models when applied to machine translation evaluation. These findings demonstrate that evaluation quality depends primarily on how well a method is matched to the specific task, rather than on model complexity alone.

At the same time, several open questions remain. In particular, it is unclear why intrinsic and extrinsic evaluation lead to contradictory conclusions about embedding quality. A deeper understanding of this mismatch would require a more precise characterization of what aspects of meaning are captured by human evaluation and how these relate to current intrinsic benchmarks. We leave this as future work: designing intrinsic evaluation methods that better reflect the requirements of downstream tasks such as machine translation evaluation.

Another promising direction for future work is the deeper exploration of sentence representations with interpretable geometric properties, as observed in our analysis of embedding spaces. Understanding what training conditions lead to these structured spaces---and how to mitigate their failures in areas like logical negation---may contribute to better generalization and more interpretable models.
---

%%% Bibliography
%%% Bibliography (literature used as a source)
%%%
%%% We employ bibTeX to construct the bibliography. It processes
%%% citations in the text (e.g., the \cite{...} macro) and looks up
%%% relevant entries in the bibliography.bib file.
%%%
%%% The \bibliographystyle command selects, which style will be used
%%% for references from the text. The argument in curly brackets is
%%% the name of the corresponding style file (*.bst). Both styles
%%% mentioned in this template are included in LaTeX distributions.

\bibliographystyle{plainnat}    %% Author (year)
% \bibliographystyle{unsrt}     %% [number]

\renewcommand{\bibname}{Bibliography}

%%% Generate the bibliography. Beware that if you cited no works,
%%% the empty list will be omitted completely.

\bibliography{bibliography}

%%% If case you prefer to write the bibliography manually (without bibTeX),
%%% you can use the following. Please follow the ISO 690 standard and
%%% citation conventions of your field of research.

% \begin{thebibliography}{99}
%
% \bibitem{lamport94}
%   {\sc Lamport,} Leslie.
%   \emph{\LaTeX: A Document Preparation System}.
%   2nd edition.
%   Massachusetts: Addison Wesley, 1994.
%   ISBN 0-201-52983-1.
%
% \end{thebibliography}


\renewcommand{\bibname}{Bibliography}
\begin{thebibliography}{13}
\small

\bibitem[Barančíková and Bojar(2019)]{barancikova-bojar-linear}
Petra Barančíková and Ondřej Bojar.
\newblock In search for linear relations in sentence embedding spaces.
\newblock In \emph{Proceedings of the 19th Conference ITAT 2019: Slovenskočeský NLP Workshop}, pages 125--132, Košice, Slovakia, 2019.

\bibitem[Barančíková and Bojar(2020a)]{costra10}
Petra Barančíková and Ondřej Bojar.
\newblock \textsc{Costra} 1.0: A dataset of complex sentence transformations.
\newblock In \emph{Proceedings of the Twelfth Language Resources and Evaluation Conference}, pages 3535--3541, Marseille, France, 2020.

\bibitem[Barančíková and Bojar(2020b)]{costra11}
Petra Barančíková and Ondřej Bojar.
\newblock \textsc{Costra} 1.1: An inquiry into geometric properties of sentence spaces.
\newblock In \emph{Proceedings of Text, Speech, and Dialogue}, pages 135--143, Cham, Switzerland, 2020.

\bibitem[Barančíková and Bojar(2025)]{barancikova-bojar-2025-intrinsic}
Petra Barančíková and Ondřej Bojar.
\newblock Intrinsic vs. extrinsic evaluation of Czech sentence embeddings: Semantic relevance doesn't help with MT evaluation.
\newblock In \emph{Proceedings of Machine Translation Summit XX: Volume 1}, pages 265--275, Geneva, Switzerland, 2025.

\bibitem[Barančíková et~al.(2014)]{barancikova:2014}
Petra Barančíková, Rudolf Rosa, and Aleš Tamchyna.
\newblock Improving evaluation of English--Czech MT through paraphrasing.
\newblock In \emph{Proceedings of LREC 2014}, pages 3402--3407, Reykjavík, Iceland, 2014.

\bibitem[Barančíková and Rosa(2015)]{barancikova-rosa-2015-targeted}
Petra Barančíková and Rudolf Rosa.
\newblock Targeted paraphrasing on deep syntactic layer for MT evaluation.
\newblock In \emph{Proceedings of Depling 2015}, pages 20--27, Uppsala, Sweden, 2015.

\bibitem[Bojar et~al.(2013)]{bojar-scratching}
Ondřej Bojar, Christian Buck, Chris Callison-Burch, Barry Haddow, Philipp Koehn, Christof Monz, Matt Post, Radu Soricut, and Lucia Specia.
\newblock Scratching the surface of possible translations.
\newblock In \emph{Proceedings of WMT 2013}, pages 432--437, Sofia, Bulgaria, 2013.


\bibitem[Kauchak and Barzilay(2006)]{kauchak}
David Kauchak and Regina Barzilay.
\newblock Paraphrasing for automatic evaluation.
\newblock In \emph{Proceedings of HLT-NAACL 2006}, pages 455--462, New York City, 2006.

\bibitem[Denkowski and Lavie(2014)]{denkowski-lavie-2014-meteor}
Michael Denkowski and Alon Lavie.
\newblock Meteor Universal: Language specific translation evaluation for any target language.
\newblock In \emph{Proceedings of WMT 2014}, pages 376--380, Baltimore, Maryland, 2014.

\bibitem[Marie et~al.(2021)]{marie-etal-2021-scientific}
Benjamin Marie, Atsushi Fujita, and Raphael Rubino.
\newblock Scientific credibility of machine translation research: A meta-evaluation of 769 papers.
\newblock In \emph{Proceedings of ACL-IJCNLP 2021}, pages 7297--7306, Online, 2021.

\bibitem[Mikolov et~al.(2013)]{word2vec}
Tomas Mikolov, Kai Chen, Greg Corrado, and Jeffrey Dean.
\newblock Efficient estimation of word representations in vector space.
\newblock In \emph{CoRR}, abs/1301.3781, 2013.

\bibitem[Papineni et~al.(2002)]{papineni-etal-2002-bleu}
Kishore Papineni, Salim Roukos, Todd Ward, and Wei-Jing Zhu.
\newblock \bleu: a method for automatic evaluation of machine translation.
\newblock In \emph{Proceedings of ACL 2002}, pages 311--318, Philadelphia, Pennsylvania, 2002.

\bibitem[Rei et~al.(2020)]{rei-etal-2020-comet}
Ricardo Rei, Craig Stewart, Ana C. Farinha, and Alon Lavie.
\newblock \comet: A neural framework for MT evaluation.
\newblock In \emph{Proceedings of EMNLP 2020}, pages 2685--2702, Online, 2020.

\bibitem[Sellam et~al.(2020)]{sellam2020bleurt}
Thibault Sellam, Dipanjan Das, and Ankur Parikh.
\newblock \bleurt: Learning robust metrics for text generation.
\newblock In \emph{Proceedings of ACL 2020}, pages 7881--7892, Online, 2020.

\end{thebibliography}

\chapwithtoc{List of Publications}
\begingroup
\small
\setlength{\parskip}{0pt}

\textbf{Publications Directly Related to the Thesis}

\begin{enumerate}
  \setlength{\itemsep}{0pt}
  \setlength{\parsep}{0pt}
    \item Petra Barančíková and Ondřej Bojar. Intrinsic vs. Extrinsic Evaluation of Czech Sentence Embeddings: Semantic Relevance Doesn't Help with MT Evaluation. In \emph{Proceedings of Machine Translation Summit XX: Volume 1}, pages 265--275, Geneva, Switzerland, August 2025. European Association for Machine Translation.

    \item Jiří Bednař, Jakub Náplava, Petra Barančíková, and Ondřej Lisický. Some Like It Small: Czech Semantic Embedding Models for Industry Applications. In \emph{Proceedings of the AAAI Conference on Artificial Intelligence}, 38(21), pages 22734--22742, 2024.

    \item Petra Barančíková and Ondřej Bojar. Costra 1.1: An Inquiry into Geometric Properties of Sentence Spaces. In \emph{Proceedings of the 23rd International Conference on Text, Speech and Dialogue}, pages 135--143, Cham, Switzerland, September 2020. Springer.

    \item Petra Barančíková and Ondřej Bojar. COSTRA 1.0: A Dataset of Complex Sentence Transformations. In \emph{Proceedings of the Twelfth Language Resources and Evaluation Conference}, pages 3535--3541, Marseille, France, May 2020. European Language Resources Association.

    \item Petra Barančíková and Ondřej Bojar. In Search for Linear Relations in Sentence Embedding Spaces. In \emph{Proceedings of the 19th Conference ITAT 2019: Slovenskočeský NLP Workshop (SloNLP 2019)}, pages 125--132, Košice, Slovakia, September 2019. CreateSpace Independent Publishing Platform.

    \item Aleš Tamchyna and Petra Barančíková. Manual and Automatic Paraphrases for MT Evaluation. In \emph{Proceedings of the Tenth International Conference on Language Resources and Evaluation (LREC'16)}, pages 3543--3548, Portorož, Slovenia, May 2016. European Language Resources Association.

    \item Petra Barančíková and Rudolf Rosa. Targeted Paraphrasing on Deep Syntactic Layer for MT Evaluation. In \emph{Proceedings of the Third International Conference on Dependency Linguistics (Depling 2015)}, pages 20--27, Uppsala, Sweden, August 2015. Uppsala University.

    \item Petra Barančíková. Parmesan: Meteor without Paraphrases with Paraphrased References. In \emph{Proceedings of the Ninth Workshop on Statistical Machine Translation}, pages 355--361, Baltimore, Maryland, USA, June 2014. Association for Computational Linguistics.

    \item Petra Barančíková, Rudolf Rosa, and Aleš Tamchyna. Improving Evaluation of English--Czech MT through Paraphrasing. In \emph{Proceedings of the Ninth International Conference on Language Resources and Evaluation (LREC'14)}, pages 3402--3407, Reykjavík, Iceland, May 2014. European Language Resources Association.

    \item Petra Barančíková and Aleš Tamchyna. Machine Translation within One Language as a Paraphrasing Technique. In \emph{Proceedings of the 14th Conference ITAT 2014}, pages 1--6, Prague, Czechia, September 2014. Institute of Computer Science AS CR.

    \item Petra Barančíková. Lexical Paraphrasing for Improvement of MT Evaluation. In \emph{WDS'13 Proceedings of Contributed Papers}, pages 7--11, Prague, Czechia, 2013. Matfyzpress.
\end{enumerate}

\textbf{Other Publications}

\begin{enumerate}
  \setlength{\itemsep}{0pt}
  \setlength{\parsep}{0pt}
    \item Václava Kettnerová, Markéta Lopatková, Anna Vernerová, and Petra Barančíková. Towards a Semi-Automatic Detection of Reflexive and Reciprocal Constructions and Their Representation in a Valency Lexicon. In \emph{Proceedings of the Twelfth Language Resources and Evaluation Conference}, pages 3136--3144, Marseille, France, May 2020. European Language Resources Association.

    \item Václava Kettnerová, Markéta Lopatková, Eduard Bejček, and Petra Barančíková. Enriching VALLEX with Light Verbs: From Theory to Data and Back Again. \emph{The Prague Bulletin of Mathematical Linguistics, 111}, pages 29--56, 2018.

    \item Petra Barančíková and Václava Kettnerová. Paraphrases of Verbal Multiword Expressions: The Case of Czech Light Verbs and Idioms. In \emph{Multiword Expressions at Length and in Depth: Extended Papers from the MWE 2017 Workshop}, pages 35--59, Berlin, Germany, 2018. Language Science Press.

    \item Petra Barančíková and Václava Kettnerová. ParaDi: Dictionary of Paraphrases of Czech Complex Predicates with Light Verbs. In \emph{Proceedings of the 13th Workshop on Multiword Expressions (MWE 2017)}, pages 1--10, Valencia, Spain, April 2017. Association for Computational Linguistics.

    \item Václava Kettnerová, Petra Barančíková, and Markéta Lopatková. Lexicographic Description of Complex Predicates in Czech: Between Lexicon and Grammar. In \emph{Proceedings of the XVII EURALEX International Congress: Lexicography and Linguistic Diversity}, pages 893--904, Tbilisi, Georgia, September 2016. Tbilisi University Press.
\end{enumerate}
  \endgroup

\end{document}